\documentclass{article}
\usepackage[T2A]{fontenc}
\usepackage[utf8]{inputenc}   
\usepackage[english, russian]{babel}

% Set page size and margins
% Replace `letterpaper' with`a4paper' for UK/EU standard size
\usepackage[a4paper,top=2cm,bottom=2cm,left=2cm,right=2cm,marginparwidth=1.75cm]{geometry}

\usepackage{amsmath}
\usepackage{graphicx}
\usepackage[colorlinks=true, allcolors=blue]{hyperref}
\usepackage{amsfonts}
\usepackage{amssymb}
% \usepackage[left=1cm,right=1cm,top=1cm,bottom=1cm]{geometry}
\usepackage{hyperref}
\usepackage{seqsplit}
\usepackage[dvipsnames]{xcolor}
\usepackage{enumitem}
\usepackage{algorithm}
\usepackage{algpseudocode}
\usepackage{algorithmicx}
\usepackage{mathalfa}
\usepackage{mathrsfs}
\usepackage{dsfont}
\usepackage{caption,subcaption}
\usepackage{wrapfig}
\usepackage[stable]{footmisc}
\usepackage{indentfirst}
\usepackage{rotating}
\usepackage{pdflscape}
\usepackage{listings}
\usepackage{amsmath}
\usepackage{graphicx}
\usepackage[colorlinks=true, allcolors=blue]{hyperref}
\usepackage{amsfonts}
\usepackage{amssymb}
% \usepackage[left=1cm,right=1cm,top=1cm,bottom=1cm]{geometry}
\usepackage{hyperref}
\usepackage{seqsplit}
\usepackage[dvipsnames]{xcolor}
\usepackage{enumitem}
\usepackage{algorithm}
\usepackage{algpseudocode}
\usepackage{algorithmicx}
\usepackage{mathalfa}
\usepackage{mathrsfs}
\usepackage{dsfont}
\usepackage{caption,subcaption}
\usepackage{wrapfig}
\usepackage[stable]{footmisc}
\usepackage{indentfirst}
\usepackage{rotating}
\usepackage{pdflscape}
\usepackage{graphicx}

\usepackage{minted}

\usepackage{MnSymbol,wasysym}
\usepackage{tikz}

\usepackage{minted}

\usepackage{MnSymbol,wasysym}

\title{Домашнее задание 11}
\author{Лиза Фоменко}
\date{}

\begin{document}

\maketitle


\section*{\textbf{Задание 1}} \textbf{\textit{Рассмотрим алгоритм Дейкстры с модификацией, в которой он сохраняет для каждой вершины предыдущую на кратчайшем пути в массиве предков $\pi[v]$. Докажите, что граф, состоящий из вершин исходного и ребер вида $\pi[v] \rightarrow v$, является деревом.}}

\textbf{Дерево} - это связный ациклический граф.

Некоторые замечания: 1) если существует вершина, не достижимая из вершины-начала алгоритма, то для нее нет значения в $\pi$, так как расстояния не существует. Согласно Вашей постановке задачи такая вершина будет просто болтаться в графе, то есть он связным не будет. 2) опять же судя по Вашей постановке непонятно, как мы проводим ребра $\pi[v] -> v$.  Для доказательства дерева эти ребра очевидно должны быть ненправленными, так как иначе мы получаем DAG.

Поэтому для доказательства прописываю, что 1) в графе нет вершин, недостижимых из вершины-начала алгоритма; 2) итоговый граф по построению неориентированный

Тогда нам нужно доказать, что граф, который мы получим, 1) связный; 2) ациклический

Докажем

1) что граф связный, то есть существует путь от любой вершины до любой другой.

по построению ребер $\pi[v] - v$ мы соединяем вершину с ее предшественником в минимальном пути $s -> v$. Алгоритм Дейкстры находит минимальный путь от $s$ до любой (достижимой) вершины $v$, то есть существует путь из $s$ до любой вершины графа. Все эти ребра входят в $\pi[v]$ по построению. Так как граф мы строим неориентированным, то также существует путь из любой вершины $v$ в $s$, из этих фактов следует, что существует путь $v - > u, u - > v$ для любой пары вершин.

2) граф ациклический. Докажем от противного, пусть в нашем получившимся графе получился цикла, то есть пускай есть цикл $v_i, ... v_k, v_i$. Это значит, что: $\pi[v_{i+1}] = v_i, ... \pi[v_{k}] = v_i$. Пусть существует путь $s -> v_i$, тогда по алгоритму Дейкстры и правилу релаксации: $d(s, v_k) \geq d(s, v_{k-1}) \geq d(s, v_{i+1}) \geq d(s, v_i)$. Тогда, по построению $\pi$, мы релаксировали ребро $v_k - v_i$, тогда верно, что $d(s, v_i) = d(s, v_k) + d(v_k, v_i)$. Но ведь мы уже доказали, что $d(s, v_k) \geq d(s, v_{k-1}) \geq d(s, v_{i+1}) \geq d(s, v_i)$, а веса всех ребер положительны, значит эти условия не могут быть выполнены одновремено. Значит, в графе не может быть циклов.

Тогда мы получили, что полученный граф - связный и ациклический, то есть является деревом.

\bigskip

\section*{\textbf{Задание 2}} \textbf{\textit{Предложите $O(|V | + |E|)$ алгоритм поиска кратчайших расстояний от данной вершины $s$ до всех остальных в графе с весами ребер $0$ или $1$. Докажите его корректность и оцените асимптотику.}}.

Насколько я помню, алгоритм называется 0-1 BFS. он создан специально для таких графов, с ребрами весом только 0 и 1, и, как указано в названии, работает через BFS за $O(|V| + |E|)$ соответственно.

Он работает на основе структуры dequeue.

\begin{verbatim}

def 01_BFS(graph, dist, start_v):
    
    dist[start_v] = 0          # от start до start 0
    dq = deque()               # это dequeue

    # кладем стартовую вершину в начало очереди
    dq.appendleft(vertex_source)

    # пока в dequeue находятся к-л вершины
    while dq:
        node = dq.popleft()        # достаем первую в очереди вершину
      
        # проверяем на релаксацию
        for edge in graph[node]:
            childNode, weight = edge     # сосед + вес ребра
          
            # аналог Дейкстры - релаксация, если нужно
            if dist[childNode] > dist[node] + weight:
                dist[childNode] = dist[node] + weight
              
                # если вес ребра = 1 (максимальный из возможных), то кладем В КОНЕЦ очереди
                if weight:
                    dq.append(childNode)
              
                # если вес ребра = 0 (минимальный из возможных), то кладем в НАЧАЛО очереди
                else:
                    dq.appendleft(childNode)
\end{verbatim}

Когда мы делаем обычный BFS, то обходим вершины по уровням, то есть в любой момент времени у нас в очереди находятся вершины либо только одного уровня, либо только двух соседних слоев. В случае dequeue в очереди одновременно находятся только элементы с двумя разыми значениями current\_weight.

В нашем случае мы используем dequeue (как и в обычной очереди добавление/вытаскивание элемента за O(1)). Рассмотрим начальную вершину. Рассмотрим всех ее соседей по ребрам с весами 0 или 1. Если вес ребра = 0, кладем вершину в начало очереди (appendleft), если = 1, то кладем в конец очереди (append). Что мы получаем на первом шаге? Сначала идут только вершины на расстоянии 0, а только затем - на расстоянии 1. Наша dequeue отсортирована. И сортировка сохраняется а каждом шаге.

Остается доказать, что, когда мы достаем вершину из очереди, для нее уже вычислено наименьшее расстояние от start.

Пусть мы вытащили вершину $v$ из начала очереди. Она имеет наименьшее расстояние $d_{curr}(s, v)$ от s из всех, находящихся в очереди (за ней могут быть вершины с такой же длиной пути, но все еще минимальное в очереди, см выше). Тогда предположим, что существует вершина $u$ с расстоянием $d(s, u)$ такая, что $d_{curr}(s, v) > d(s, u) + d(u, v)$. Мы знаем, что $d(u, v) \in \{0, 1\}$. Но мы знаем, что вершина $u$ лежит в очереди после вершины $v$, значит, $d(s, u) \geq d_{curr}(s, v)$. Но тогда мы получаем противоречие, а значит, $d(s,v)$ является минимальным. Тогда, когда мы достаем вершину из очереди, для нее уже вычислено минимальное расстояние.

Тогда перейдем к ассимптотике. Каждое ребро релаксируется не более 1 раза, поэтому каждая вершина заносится в очередь не более раза. То есть мы по одному разу достаем вершину из очереди, для каждой вершины смотрим все ее ребра (всего для всех вершин |E|). Релаксация ребра и занесение в очередь производится за O(1). Тогда общая ассимптотика, как и требовалось, $O(|V| + |E|)$

\bigskip


\section*{\textbf{Задание 3}} \textbf{\textit{Изучите топологическую сортировку (она делается на основе обхода в глубину). Рекомендуется использовать учебники Дасгупты и Кормена. Сделайте топологическую сортировку графа}}

DFS был проведен в алфавитом порядке

\begin{figure}[H]
     \centering
     \begin{subfigure}[b]{0.95\textwidth}
         \centering
         \includegraphics[width=\textwidth]{algos/toposort.png}
         \caption{Toposort (inspired by Kormen)}
     \end{subfigure}
     \hfill
\end{figure}


\bigskip

\section*{\textbf{Задание 4}} \textbf{\textit{В графе может быть несколько кратчайших путей между какими-то вершинами. Постройте линейный по времени алгоритм, находящий количество вершин, которые лежат \textbf{хотя бы на одном} кратчайшем пути из $s$ в $t$ в неориентированном графе с единичными весами на рёбрах.}}

Так как весь ребер единичный, то мы можем считать длину пути просто за число шагов в BFS. Если бы мы шли за один BFS, то пришлось бы для каждой вершины хранить массив вершин, что не хочется. Так как нам в принципе не важно расположение путей, то заметим следующее: $d(s, v) + s(v, t) = d(s, t)$. 

Проведем BFS с поиском кратчайшего пути (как только мы впервые дошли до вершины -- нашли наименьший путь). Начинаем из s, для каждой вершины найдем $d(s, v)$, а также вычислим $d(s, t)$. Затем второй BFS проведем с началом в $t$, получим $d(t, v)$. Так как граф неориентированный, то $d(t, v) = d(v,t)$. Оба BFS работают за $O(|E| + |V|)$. Затем мы можем для каждой вершины проверить факт $d(s, v) + s(v, t) = d(s, t)$, ведь мы за два BFS посчитали все три расстояния из выражения. Если для к-л вершины выполняется равенство, то она лежит на каком-то минимальном пути. Такая проверка занимает $O(|V|)$, Total  $O(|E| + |V|)$

\bigskip

\section*{\textbf{Задание 5}} \textbf{\textit{Примените алгоритм Дейкстры к графу для поиска кратчайших путей от вершины $e$ до всех остальных}}

Сделано по алгоритму Дейктрсы, но пишу не все шаги, а только те, которые в итоге приводят к наименьшим путям для вершин (соответствющие ребра подписаны iter N)

\begin{figure}[H]
     \centering
     \begin{subfigure}[b]{0.95\textwidth}
         \centering
         \includegraphics[width=\textwidth]{algos/dijkstra.png}
         \caption{Dijkstra's algorithm}
     \end{subfigure}
     \hfill
\end{figure}



\bigskip

\section*{\textbf{Задание 6}} \textbf{\textit{Профессор О. П. Рометчивый предлагает следующий способ нахождения кратчайшего пути из $s$ в $t$ в ориентированном графе, содержащем рёбра отрицательного веса. Прибавим достаточно большую константу к весам всех рёбер и сделаем все веса положительными, после чего воспользуемся алгоритмом Дейкстры. Корректен ли такой подход? Если да, то докажите это, если нет — приведите контрпример.}}

Это просто неверно. Типа да, такая ситуация спасает отрицательные ребра, но граф, который мы получаем, совсем другой. Но в случаях, когда мы добавляем очень большую константу C, которая, условно, в 100 раз больше максимального элемента (и перебрасывает минимальный в положительные числа при добавлении), то что мы можем сказать? Пусть у нас был какой-то путь веса $w_1 + w_2 + w_3... + w_k$. Прибавив к каждому ребру константу, мы получим, что новый вес пути $w_1 + w_2 + w_3... + w_k + k*C$. Если C достаточно большая, то $w_i$ могут выглядеть не более чем погрешность на ее фоне. Но тогда мы получим, что вес пути прямо пропорционален его длине, что в корне неверно, смотрим контрпример.

\begin{figure}[H]
     \centering
     \begin{subfigure}[b]{0.95\textwidth}
         \centering
         \includegraphics[width=\textwidth]{algos/task6.png}
     \end{subfigure}
     \hfill
\end{figure}

Этот рисунок довольно тривиален. Мы видим, что в реальности (левый рисунок) самый короткий путь - путь2 длины 6. Но прибавив даже минимальную константу, которая передвет отрицательное ребро в положительное, мы получаем граф на рисунке 2 - в нем минимальным становится путь2. Это в принципе неверно

\bigskip

\end{document}